\section{Chương~\ref{sec:muscles}. Đầu ra tới cơ}

Cấu tạo của đầu ra (đơn vị vận động, hệ số an toàn của khớp thần kinh, bật theo kích thước, cặp cơ, vòng phản xạ giãn và bộ tạo nhịp) đã được đo trực tiếp; điều kiện ổn định của vòng có trễ là một định lý; mô hình bên trong của cơ thể là giả thuyết có cơ sở vững; các con số trên hình được tính theo mô hình của sách.

{\footnotesize
\setlength{\tabcolsep}{3pt}\renewcommand{\arraystretch}{1.15}
\begin{longtable}{>{\raggedright}p{0.5cm}>{\raggedright}p{4.0cm}>{\raggedright}p{5.6cm}>{\raggedright}p{2.3cm}>{\raggedright\arraybackslash}p{1.7cm}}
\toprule
TT & Khẳng định & Thí nghiệm và điều đã đo & Nguồn & Trạng thái \\
\midrule\endfirsthead
\toprule
TT & Khẳng định & Thí nghiệm và điều đã đo & Nguồn & Trạng thái \\
\midrule\endhead
1 & Đơn vị vận động: một nơ-ron \nt{ACET} điều khiển một nhóm sợi, từ vài trăm đến vài nghìn; ở cơ cần tinh chỉnh, đơn vị nhỏ hơn & Cơ người, đếm sợi trục vận động và sợi cơ: trung bình khoảng $340$ sợi trên một nơ-ron ở cơ gian cốt bàn tay, khoảng $410$ ở cơ cánh tay quay (brachioradialis), khoảng $1930$ ở cơ bụng chân trong. Chương không nêu con số chính xác cho những đơn vị nhỏ nhất. & \cite{Feinstein1955mu} & đã đo \\
2 & Khớp thần kinh trên cơ giải phóng lượng chất dẫn truyền gấp vài lần mức cần để sợi cơ đáp ứng & Tổng quan các phép đo ở khớp thần kinh--cơ: hệ số an toàn (tỷ số giữa lượng chất dẫn truyền được giải phóng và lượng ngưỡng) ở động vật có vú trưởng thành thường là $3$--$5$; ở người, hệ số an toàn dựa vào các nếp gấp của màng bên nhận nhiều hơn là vào lượng giải phóng. & \cite{WoodSlater2001} & đã đo \\
3 & Cú giật~\eqref{eq:twitch} tăng lên trong thời gian $T_c$ cỡ $50\ms$ & Từng đơn vị vận động bàn tay người khi gắng sức tự chủ, lực được lấy trung bình theo xung của đơn vị: thời gian tăng $30$--$100\ms$, ở $80\,\%$ đơn vị dưới $70\ms$. Hàm $(t/T_c)e^{1-t/T_c}$ là xấp xỉ lấy từ mô hình quần thể đơn vị. & \cite{MilnerBrown1973rec, Fuglevand1993} & đã đo \\
4 & Tổng các cú giật~\eqref{eq:forcesum}: với $5$, $15$ và $40$ xung/s, lực là $0.2$--$1.1F_1$, $1.8$--$2.2F_1$ và khoảng $5.4F_1$ (hình~\ref{fig:tetanus}) & Tính theo \texttt{models/\allowbreak muscle\_\allowbreak models.py}, $T_c=50\ms$: $0.21$--$1.10$, $1.76$--$2.19$ và $5.28$--$5.49\,F_1$ trong $0.3\,\text{s}$ cuối. Phép cộng tuyến tính là đơn giản hóa: mô hình không có sự bão hòa của cơ thật. & \texttt{models/\allowbreak muscle\_\allowbreak models.py} & mô hình \\
5 & Các đơn vị được bật theo kích thước; đơn vị yếu nhất yếu hơn đơn vị mạnh nhất cả trăm lần & Rễ trước tủy sống của mèo: khi đầu vào phản xạ tăng dần, những nơ-ron có xung nhỏ trong rễ phóng điện trước, tức là nơ-ron nhỏ. Cơ gian cốt bàn tay người: lực giật của các đơn vị từ $0.1$ đến $10$~g và tăng gần tuyến tính theo mức lực tại đó đơn vị được bật. & \cite{Henneman1957, MilnerBrown1973rec} & đã đo \\
6 & Bước lực tỉ lệ với lực, $\Delta F/F\approx q-1$~\eqref{eq:sizeprinciple}: dấu phẩy động & Suy ra trong chương từ cấp số nhân của các lực. Phù hợp với thực nghiệm: khi gắng sức lớn, cùng một mức tăng lực chỉ cần bật thêm ít đơn vị mới hơn nhiều. & suy ra trong chương; \cite{MilnerBrown1973rec} & lý thuyết \\
7 & Độ tản của lực tăng tỉ lệ với chính lực & Gắng sức đẳng trường tự chủ của cơ ngón cái: độ lệch chuẩn của lực tăng tuyến tính theo lực trung bình; với cùng mức lực tạo bằng kích thích điện, độ tản không đổi. Mô hình quần thể: tăng tuyến tính là do bật đơn vị theo thứ tự lực. & \cite{Jones2002sdn} & đã đo \\
8 & Sự mượt mà của cử động là hệ quả của nhiễu phụ thuộc tín hiệu & Mô hình: quỹ đạo cực tiểu hóa độ tản ở điểm cuối dưới loại nhiễu này tái hiện các quỹ đạo đã đo của mắt và tay. Chưa có thí nghiệm trực tiếp tách nguyên nhân này khỏi các nguyên nhân khác. & \cite{HarrisWolpert1998} & giả thuyết \\
9 & Hiệu lực của cặp cơ quyết định mô-men, tổng quyết định độ cứng~\eqref{eq:antagonists} & Lý thuyết điều khiển trở kháng bằng đồng co. Cánh tay người: độ cứng của mỗi khớp, đo bằng những cú đẩy nhỏ, xấp xỉ tỉ lệ với mô-men tại khớp đó; các tỷ lệ đồng co khác nhau thay đổi hình dạng và hướng của elip độ cứng ở bàn tay. & \cite{Hogan1984, GomiOsu1998stiff} & đã đo \\
10 & Cảm biến độ giãn kích thích nơ-ron \nt{ACET} của nó và qua nơ-ron trung gian ức chế tắt cơ đối vận (hình~\ref{fig:antagonists}) & Tủy sống mèo: một nơ-ron trung gian đơn lẻ, được kích thích bởi sợi của cảm biến độ giãn, tạo điện thế ức chế qua một khớp thần kinh ở các nơ-ron \nt{ACET} của cơ đối vận, độ trễ qua khớp thần kinh $0.28$--$0.42\ms$. Chất dẫn truyền của nơ-ron trung gian (glycine) không được xác định trong thí nghiệm này. & \cite{Jankowska1972Ia} & đã đo \\
11 & Độ trễ của vòng: tủy sống $25$--$30\ms$, qua vỏ não dài hơn một rưỡi đến ba lần, qua mắt $100$--$200\ms$ & Cánh tay người, đẩy vào khớp: đáp ứng ngắn của cơ sau $20$--$45\ms$, đáp ứng dài $45$--$105\ms$, đáp ứng tự chủ $120$--$180\ms$. Dịch vị trí nhìn thấy của ngón tay trong lúc cử động: hiệu chỉnh trung bình sau $160\ms$. & \cite{Pruszynski2008rapid, SaundersKnill2003vis} & đã đo \\
12 & $\dd x/\dd t=-Gx(t-d)$ ổn định khi $Gd<\pi/2$~\eqref{eq:delaystable}; tại ranh giới, tần số là $1/(4d)$ & Định lý về nghiệm của phương trình có trễ. Kiểm tra theo \texttt{models/\allowbreak muscle\_\allowbreak models.py}, $d=30\ms$: tới $3\,\text{s}$, biên độ là $3\cdot10^{-6}$ khi $G=40$, $0.13$ khi $G=50$, $38$ khi $G=55$ mỗi giây (ranh giới $52$). & \cite{Hayes1950} & lý thuyết \\
13 & Run sinh lý $8$--$12$~Hz; vòng phản xạ giãn là một trong các nguồn của nó & Tổng quan các phép đo: run nhẹ trong dải quanh $10$~Hz có ở người khỏe mạnh; nguồn gốc của nó hỗn hợp (dao động trung ương, đặc tính phóng điện của đơn vị, cộng hưởng cơ học và cộng hưởng của vòng phản xạ), và trong những điều kiện khác nhau thì yếu tố khác nhau chiếm ưu thế. & \cite{McAuleyMarsden2000} & giả thuyết \\
14 & Não dự đoán kết quả của lệnh theo mô hình bên trong của cơ thể & Người cầm vật nặng bằng cách kẹp ngón và cử động tay: với tải quán tính, nhớt và đàn hồi, lực kẹp thay đổi đồng thời với tải chứ không trễ theo vòng phản hồi, tức là tải đã được dự đoán. Bài tổng quan tập hợp các thí nghiệm như vậy; bản thân mô hình chưa được quan sát trực tiếp trong nơ-ron. & \cite{FlanaganWing1997grip, WolpertGhahramani2000} & giả thuyết \\
15 & Nhịp đi, thở, nhai do các mạng cục bộ sinh ra khi đầu vào không đổi & Tổng quan thí nghiệm: hệ thần kinh bị cô lập (tủy sống, hạch) tạo ra mẫu vận động có nhịp mà không cần đầu vào có nhịp và không cần phản hồi từ cơ. & \cite{MarderBucher2001} & đã đo \\
16 & Ức chế lẫn nhau kèm sự mỏi~\eqref{eq:halfcentre} cho nhịp với chu kỳ $0.84\,\text{s}\approx2\tau_a$ (hình~\ref{fig:halfcentre}) & Định lý: mạng ức chế lẫn nhau có thích nghi, không có trạng thái ổn định khi đầu vào không đổi, nhất thiết phải dao động. Tính theo \texttt{models/\allowbreak muscle\_\allowbreak models.py} ($I=1$, $\beta=2$, $b=2$, $\tau=20\ms$, $\tau_a=0.4\,\text{s}$): chu kỳ $0.84\,\text{s}$. Việc bộ tạo nhịp thật được cấu tạo đúng như vậy chưa được chứng minh. & \cite{Matsuoka1985osc} & mô hình \\
\bottomrule
\end{longtable}}
